% ============================================================
% TABLE 10 — SYNTHESIS OF INVARIANCE
% Book I: Ontology of Viability
% Part I: Identity and Manifestation
% Chapter 2: Invariance
% ============================================================

\begin{table}[htbp]
    \centering
    \caption{Conceptual Synthesis of Invariance}
    \label{tab:synthesis-of-invariance}

    \renewcommand{\arraystretch}{1.4}

    \begin{tabular}{|p{0.22\textwidth}|p{0.36\textwidth}|p{0.33\textwidth}|}
        \hline
        \textbf{Conceptual Relation} &
        \textbf{Established Distinction} &
        \textbf{Result for the Concept of Invariance} \\
        \hline

        Invariance and identity &
        Identity concerns the existent as identity-bearing; invariance concerns the remaining unchanged of a specified determination under specified transformations. &
        Invariance does not constitute identity or exhaust its meaning. \\
        \hline

        Invariance and difference &
        Difference concerns determinate distinctions, whereas invariance concerns the unchanged status of a particular respect within a specified domain. &
        Difference and invariance may coexist, provided that they concern distinct respects or determinations. \\
        \hline

        Invariance and change &
        Change concerns determinate alteration; invariance concerns the remaining unchanged of a specified determination under a specified transformation. &
        An existent may change in one respect while remaining invariant in another. \\
        \hline

        Invariance and continuity &
        Continuity concerns the continuation of an existent or process; invariance concerns the unchanged status of a particular determination. &
        Invariance does not, by itself, establish continuity, and continuity does not necessarily require a specified invariant feature. \\
        \hline

        Invariance and structure &
        Structure concerns an organized configuration of determinations and relations; structural invariance concerns the preservation of a specified structural determination or relation. &
        The invariance of a structural component does not establish the invariance of the whole structure. \\
        \hline

        Invariance and knowledge &
        Ontological invariance concerns the status of what remains unchanged; epistemic recognition concerns the ways in which a knower identifies and describes that invariance. &
        Knowledge does not constitute ontological invariance, and incomplete knowledge does not establish its absence. \\
        \hline

        Forms and domains of invariance &
        Quantitative, qualitative, relational, and structural forms are analytical distinctions; domains specify the transformations under consideration. &
        These forms are not assumed to be mutually exclusive or to constitute a final classification of ontological kinds. \\
        \hline

        Scope and limits &
        The relational account identifies a subject, a domain of transformations, and a respect in which something remains unchanged. &
        The account provides a working definition without settling the ultimate ontological status of invariance. \\
        \hline

    \end{tabular}
\end{table}

\vspace{0.5em}

\noindent
\textit{General synthesis:}
Invariance is provisionally understood in this chapter as a relation between an existent, or a determinate aspect of an existent, and a specified domain of changes or transformations, such that the existent or aspect remains unchanged in a particular respect under those transformations. This account distinguishes invariance from identity, difference, change, continuity, structure, and epistemic recognition without treating these concepts as wholly unrelated.

% ============================================================
% END OF TABLE 10
% ============================================================